\chapter{Introduction}
\section{Motivation}
Supervised classification learns relationships between observations and their
labels. A label is often treated as a reliable target even when it was produced
through an imperfect annotation process. Ambiguous definitions, transcription
mistakes, incomplete evidence, and disagreement between annotators can create
incorrect targets. The survey by Fr\'{e}nay and Verleysen distinguishes label
noise from other data defects and reviews its consequences for classification
and the methods used to address it~\cite{frenay2014}.

This project focuses on numeric tabular representations. Such representations
include measured chemical properties, image-derived morphology, and flattened
pixel intensities. They share a matrix interface, but their class structure
and difficulty differ. A method that works on one should not be assumed to
work on all. Nazaretyan, Leser, and Kircher's tabular benchmark reports that
no evaluated label-error identification method dominates all studied
settings~\cite{nazaretyan2025}. The practical question is therefore not merely
whether a method can flag unusual observations, but whether its flags are
useful for improving a particular downstream learner.

Two tasks must be separated. \emph{Detection} identifies training rows whose
observed labels are suspect. \emph{Intervention} removes or changes those
labels and then retrains a classifier. A high detection score does not
guarantee that intervention improves classification. Removal can discard
informative boundary examples, and relabeling can reinforce an incorrect
prediction. This distinction motivates evaluating both tasks against separate
targets: known injected corruptions for detection, and an untouched test set
for downstream prediction.

\section{Research Questions}
The study addresses three questions from the project proposal:
\begin{enumerate}
\item \textbf{RQ1:} How does controlled training-label noise affect common
tabular classifiers across datasets, corruption mechanisms, and noise rates?
\item \textbf{RQ2:} How accurately do lightweight out-of-fold confidence,
disagreement, margin, and combined scores identify the injected errors?
\item \textbf{RQ3:} Does removing or relabeling the selected examples improve
test accuracy, macro F1, and AUROC relative to training on noisy labels?
\end{enumerate}
The combined score is a specific hypothesis: assigning weights to several
signals might produce a better ordering of suspicious examples than any
component alone. It is evaluated against its individual components under an
identical selection budget. Its implementation is not, by itself, evidence
of novelty or superiority.

\section{Contributions and Scope}
The principal contribution is a reproducible controlled benchmark connecting
label-error detection to downstream classification. It includes four datasets,
three model families, two corruption mechanisms, four corruption rates, four
random seeds, four detectors, and two cleaning actions. The study also provides
a direct component comparison, an analytical explanation of the combined
score's binary behavior, and an engineering release with original training
provenance, fitted-model examples, generated figures, and checksum verification.

The empirical contribution includes a negative result: the weighted combination
does not improve average performance over confidence alone. Reporting that
outcome is essential to answering the proposed research question. The results
support a limited practical conclusion about simple filtering under the
conditions tested, not a claim that the score is a new state-of-the-art method.

The scope is deliberately controlled. Only training labels are corrupted;
features are unchanged. The original dataset labels serve as reference labels,
not independently certified ground truth. Selection uses the known injected
noise fraction. No categorical-data pipeline, unknown-noise estimation,
instance-dependent corruption, hyperparameter search, or neural model is
evaluated. These limits delimit the conclusions rather than being hidden
implementation details.

\section{Organization}
Chapter~2 situates the work within noisy-label learning and tabular evaluation.
Chapter~3 defines the corruption, detection, and cleaning procedures.
Chapter~4 specifies datasets, estimators, metrics, and statistical aggregation.
Chapter~5 analyzes the results and failure cases. Chapter~6 concludes and
identifies follow-up experiments. The appendix documents reproducibility and
provides illustrative diagnostic curves.
